% =====================================================================
%  Расчётно-пояснительная записка к ВКР бакалавра
%  Тема: Разработка программно-аппаратного комплекса поддержки
%        мобильного приложения визуализации результатов работы
%        модификаций эволюционного алгоритма оптимизации
%  Автор: А.В. Волохов, ИУ9-82Б, МГТУ им. Н.Э. Баумана
%  Руководитель: Д.П. Посевин
%
%  Преамбула приведена в соответствие с разделом 10 Положения о ВКР:
%   - шрифт Times New Roman 14 pt (пакет tempora);
%   - полуторный межстрочный интервал;
%   - поля: левое 30 мм, правое 15 мм, верхнее 20 мм, нижнее 20 мм;
%   - абзацный отступ 1,25 см, выравнивание по ширине;
%   - номер страницы внизу по центру, верхний колонтитул пустой;
%   - рисунок: подпись снизу, «Рисунок N — Название»;
%   - таблица: подпись сверху, «Таблица N — Название»;
%   - формулы: сквозная нумерация, номер справа в круглых скобках.
%
%  Сборка: pdflatex -> biber -> pdflatex -> pdflatex
% =====================================================================

\documentclass[14pt, a4paper]{extarticle}

%%% Кодировки и язык %%%
\usepackage{cmap}
\usepackage[T2A]{fontenc}
\usepackage[utf8]{inputenc}
\usepackage[english, main=russian]{babel}

%%% Шрифт Times New Roman %%%
\usepackage{tempora}

%%% Поля страницы (Положение, п. 10.2) %%%
\usepackage{geometry}
\geometry{a4paper,
  left=30mm, right=15mm, top=20mm, bottom=20mm,
  headsep=10pt, footskip=12mm}

%%% Полуторный межстрочный интервал %%%
\usepackage{setspace}
\onehalfspacing

%%% Абзацный отступ 1,25 см и отступ у первого абзаца %%%
\usepackage{indentfirst}
\setlength{\parindent}{1.25cm}

%%% Математика %%%
\usepackage{amsmath, amssymb, amsfonts, mathtools}

%%% Таблицы %%%
\usepackage{array}
\usepackage{longtable}
\usepackage{multirow}
\usepackage{makecell}
\usepackage{tabularx}
\usepackage{booktabs}

%%% Рисунки %%%
\usepackage{graphicx}
\usepackage{float}

%%% Списки %%%
\usepackage{enumitem}
\setlist{nosep}
\renewcommand{\labelitemi}{---}

%%% Кавычки-ёлочки %%%
\usepackage{csquotes}
\DeclareQuoteStyle{russian}
  {\guillemotleft}{\guillemotright}[0.025em]
  {\quotedblbase}{\textquotedblleft}
\ExecuteQuoteOptions{style=russian}
\newcommand{\enq}[1]{\enquote{#1}}

%%% Десятичная запятая в формулах %%%
\usepackage{icomma}

%%% Подписи рисунков и таблиц (Положение, пп. 10.5, 10.6) %%%
\usepackage{caption}
\DeclareCaptionLabelSeparator{emdash}{~--~}
\captionsetup[figure]{labelsep=emdash, position=bottom, name=Рисунок,
  justification=centering, singlelinecheck=false}
\captionsetup[table]{labelsep=emdash, position=top, name=Таблица,
  justification=raggedright, singlelinecheck=false}

%%% Листинги кода (для приложений) — на пакете listings %%%
\usepackage{listings}
\usepackage{xcolor}
\definecolor{lstcomment}{rgb}{0.0,0.5,0.0}
\definecolor{lstkeyword}{rgb}{0.0,0.0,0.6}
\definecolor{lststring}{rgb}{0.58,0,0.82}
\renewcommand{\lstlistingname}{Листинг}
\lstset{
  basicstyle=\ttfamily\footnotesize,
  breaklines=true,
  breakatwhitespace=false,
  frame=single,
  numbers=left,
  numbersep=6pt,
  numberstyle=\footnotesize\color{gray},
  keywordstyle=\color{lstkeyword}\bfseries,
  commentstyle=\color{lstcomment},
  stringstyle=\color{lststring},
  showstringspaces=false,
  extendedchars=true,
  inputencoding=utf8,
  captionpos=t,
  xleftmargin=18pt,
  tabsize=4,
}

%%% Колонтитулы: номер страницы внизу по центру, верх пустой (п. 10.3) %%%
\usepackage{fancyhdr}
\pagestyle{fancy}
\fancyhf{}
\renewcommand{\headrulewidth}{0pt}
\renewcommand{\footrulewidth}{0pt}
\fancyfoot[C]{\thepage}
\fancypagestyle{plain}{%
  \fancyhf{}\fancyfoot[C]{\thepage}%
  \renewcommand{\headrulewidth}{0pt}}

%%% Заголовки разделов (Положение, п. 10.4):
%%%   «глава/раздел» не пишется, в конце заголовка точки нет,
%%%   допускается полужирный шрифт, без переносов.
\usepackage[explicit]{titlesec}
\titleformat{\section}
  {\normalfont\large\bfseries\raggedright}{\thesection\quad}{0pt}{#1}
\titleformat{\subsection}
  {\normalfont\normalsize\bfseries\raggedright}{\thesubsection\quad}{0pt}{#1}
\titleformat{\subsubsection}
  {\normalfont\normalsize\bfseries\raggedright}{\thesubsubsection\quad}{0pt}{#1}
\titlespacing*{\section}{0pt}{12pt}{8pt}
\titlespacing*{\subsection}{0pt}{10pt}{6pt}
\titlespacing*{\subsubsection}{0pt}{8pt}{6pt}

% Каждая основная часть — с новой страницы
\newcommand{\sectionbreak}{\clearpage}

%%% Ненумеруемый структурный элемент (заголовок прописными, по центру,
%%% с новой страницы, заносится в оглавление) — п. 9.3
\newcommand{\structelem}[1]{%
  \clearpage
  \phantomsection
  \addcontentsline{toc}{section}{#1}%
  \begin{center}\large\bfseries #1\end{center}
  \vspace{8pt}}

%%% Ссылки %%%
\usepackage[hidelinks]{hyperref}
\usepackage{cleveref}
\crefname{figure}{рисунок}{рисунки}
\crefname{table}{таблица}{таблицы}
\crefname{equation}{формула}{формулы}
\crefformat{equation}{(#2#1#3)}

%%% Библиография: biblatex + ГОСТ (Положение, п. 9.3 — ГОСТ 7.32) %%%
\usepackage[
  backend=biber,
  bibstyle=gost-numeric,
  citestyle=numeric-comp,
  sorting=none,
  language=russian
]{biblatex}
\addbibresource{biblio.bib}

%%% Перенос слов и борьба с висячими строками %%%
\sloppy
\clubpenalty=10000
\widowpenalty=10000

% =====================================================================
\begin{document}

% --- На время разработки страницы до введения опускаем ---
% (титульный лист, задание, календарный план — вставляются по бланкам кафедры)

% =====================================================================
%  РЕФЕРАТ (АННОТАЦИЯ) — до 1 страницы, не нумеруется
% =====================================================================
\structelem{РЕФЕРАТ}
% TODO: объект и цель работы; методы; результаты; область применения;
%       сведения об объёме (страниц, разделов, рисунков, таблиц,
%       источников, приложений). Заполняется в конце.

% =====================================================================
%  СОДЕРЖАНИЕ
% =====================================================================
\clearpage
\phantomsection
\addcontentsline{toc}{section}{СОДЕРЖАНИЕ}
\renewcommand{\contentsname}{\centerline{\large\bfseries СОДЕРЖАНИЕ}}
\tableofcontents

% =====================================================================
%  СПИСОК ОБОЗНАЧЕНИЙ И СОКРАЩЕНИЙ
% =====================================================================
\structelem{СПИСОК ОБОЗНАЧЕНИЙ И СОКРАЩЕНИЙ}

В настоящей ВКР применяют следующие сокращения и обозначения.

\begin{description}[leftmargin=!, labelwidth=2.6cm, font=\normalfont, itemsep=2pt]
  \item[BBO] --- алгоритм биогеографической оптимизации (Biogeography-Based Optimization)
  \item[GIF] --- формат обмена графическими данными (Graphics Interchange Format)
  \item[HSI] --- индекс пригодности среды обитания (Habitat Suitability Index)
  \item[HTTP] --- протокол передачи гипертекста (Hypertext Transfer Protocol)
  \item[JSON] --- текстовый формат обмена данными (JavaScript Object Notation)
  \item[PSO] --- метод роя частиц (Particle Swarm Optimization)
  \item[SSE] --- технология однонаправленной потоковой передачи событий от сервера (Server-Sent Events)
  \item[TCP] --- протокол управления передачей (Transmission Control Protocol)
  \item[UI] --- пользовательский интерфейс (User Interface)
  \item[WebSocket] --- протокол полнодуплексной связи поверх TCP
\end{description}

% =====================================================================
%  ВВЕДЕНИЕ — ~1 страница, без подразделов, не нумеруется
% =====================================================================
\structelem{ВВЕДЕНИЕ}

Многие практические задачи сводятся к поиску глобального минимума целевой
функции. Часто у такой функции, помимо глобального минимума, есть множество
локальных. Классические градиентные методы здесь малоэффективны: они сходятся
к ближайшему локальному минимуму и не умеют из него выбираться. Поэтому для
функций с множеством локальных минимумов применяют метаэвристические алгоритмы,
которые работают сразу с целой популяцией решений. Один из них --- алгоритм
биогеографической оптимизации (Biogeography-Based Optimization, далее --- BBO).

Поведение метаэвристик трудно предсказать аналитически, поэтому об их работе
чаще судят по результатам запусков. Здесь многое даёт визуализация: по ней
видно, как убывает лучшее найденное значение, как популяция распределяется в
пространстве решений и не наступает ли преждевременная сходимость. Удобно, когда
такой инструмент доступен прямо на мобильном устройстве и не требует настройки
вычислительного окружения. Этим и обусловлена актуальность работы.

Объект исследования --- модификации алгоритма биогеографической оптимизации,
а также метод роя частиц, который привлекается для сравнения. Программа
рассчитана прежде всего на студентов, изучающих методы оптимизации, и на
исследователей, которым нужно быстро сравнить поведение разных метаэвристик на
тестовых задачах.

Цель работы --- разработка программно-аппаратного комплекса поддержки мобильного
приложения визуализации результатов работы модификаций эволюционного алгоритма
оптимизации. Для её достижения решается несколько связанных между собой задач.

Вычисления выполняет серверное ядро. Оно должно поддерживать целое семейство
методов: классический BBO, его варианты с синусоидальной и смешанной моделями
миграции, их модификации с адаптивной мутацией, перезапуском при стагнации и
локальным поиском, а также метод роя частиц. Запускать алгоритмы нужно как на
стандартных тестовых функциях, так и на функции, которую пользователь задаёт
собственным выражением.

Чтобы мобильное устройство могло ставить задачи серверу и получать от него
данные, нужен протокол обмена. Он должен поддерживать постановку задачи, передачу
прогресса по итерациям в реальном времени, выдачу итогового результата, отмену
задачи и работу очереди, когда клиентов несколько.

С пользователем взаимодействует мобильное приложение. Через него настраивают
эксперимент, отправляют задачу на сервер и следят за ходом вычислений и итогом.

Главное в приложении --- наглядно показать, как работает алгоритм. Для этого
строится график сходимости в реальном времени, пространство поиска отображается
контурной картой и трёхмерной поверхностью, а динамика популяции --- анимацией.

Смысл комплекса в том, чтобы сравнивать методы, поэтому история запусков
сохраняется. Сохранённые результаты можно сопоставлять --- на совмещённых
графиках и в таблице.

Наконец, комплекс применяется по назначению: проводятся вычислительные
эксперименты, в которых реализованные методы сравниваются на тестовых функциях
разной структуры и размерности, а результаты анализируются.

% =====================================================================
%  1. ИССЛЕДОВАТЕЛЬСКАЯ ЧАСТЬ
% =====================================================================
\section{Исследовательская часть}\label{sec:research}

\subsection{Задача глобальной оптимизации и метаэвристики}
\label{subsec:opt-problem}

Задача глобальной оптимизации состоит в поиске точки, в которой целевая функция
принимает наименьшее значение:
\begin{equation}\label{eq:global-min}
  f(x) \to \min, \qquad x \in D \subset \mathbb{R}^{d},
\end{equation}
где $f$ --- целевая функция, $D$ --- допустимая область поиска, $d$ ---
размерность пространства решений. Задача максимизации сводится к минимизации
заменой $f$ на $-f$, поэтому далее везде речь идёт о минимизации.

Сложность задачи определяется рельефом целевой функции. Если у функции
единственный минимум, найти его несложно. Гораздо труднее функции с множеством
локальных минимумов: наряду с глобальным минимумом у них есть много локальных,
и отличить один от другого по поведению функции в окрестности точки невозможно.
Классические градиентные методы движутся в сторону убывания функции и поэтому
сходятся к ближайшему минимуму, не различая, локальный он или глобальный. Для
таких задач они малопригодны. К тому же градиентные методы требуют, чтобы
функция была дифференцируемой, а это условие выполняется не всегда.

Альтернативой служат метаэвристические алгоритмы. Это итеративные методы поиска,
идеи которых заимствованы из природных процессов: эволюции, поведения стай,
миграции видов. Их главная особенность --- работа не с одной точкой, а с целой
популяцией решений. За счёт этого алгоритм одновременно просматривает несколько
областей пространства поиска, и вероятность застрять в локальном минимуме
снижается. Разнообразие популяции поддерживается случайными механизмами:
мутацией, обменом информацией между решениями, отбором. Метаэвристики не требуют
дифференцируемости функции и не используют производные, поэтому подходят для
широкого круга задач --- в том числе для таких, где функция задана лишь
алгоритмом её вычисления. Платят за это отсутствием гарантий: найденное решение
может оказаться не глобальным минимумом, а лишь близким к нему. К этому классу
относятся, в частности, генетические алгоритмы, метод роя частиц (Particle Swarm
Optimization, далее --- PSO) и рассматриваемый в работе алгоритм
биогеографической оптимизации.

Чтобы сравнивать алгоритмы между собой, используют наборы тестовых функций
с заранее известными минимумами. Функции подбирают так, чтобы их рельеф был
разным и проверял разные стороны поведения метода. В работе используются четыре
такие функции.

Функция сферы --- самая простая, с единственным минимумом:
\begin{equation}\label{eq:sphere}
  f_{\text{sph}}(x) = \sum_{i=1}^{d} x_i^{2}.
\end{equation}
Её глобальный минимум равен нулю и достигается в начале координат. На ней удобно
оценивать скорость сходимости алгоритма, поскольку локальных ловушек нет.

У функции Растригина минимумов очень много:
\begin{equation}\label{eq:rastrigin}
  f_{\text{ras}}(x) = A d + \sum_{i=1}^{d}
    \bigl(x_i^{2} - A \cos(2\pi x_i)\bigr), \qquad A = 10.
\end{equation}
Косинусное слагаемое создаёт большое число локальных минимумов, расположенных
правильной решёткой. Глобальный минимум, равный нулю, находится в начале
координат. Эта функция проверяет, способен ли алгоритм не застревать в локальных
минимумах.

У функции Экли один глобальный минимум, окружённый почти плоской областью
с мелкими неровностями:
\begin{equation}\label{eq:ackley}
  f_{\text{ack}}(x) = -a \exp\!\left(-b \sqrt{\frac{1}{d}\sum_{i=1}^{d} x_i^{2}}\right)
  - \exp\!\left(\frac{1}{d}\sum_{i=1}^{d} \cos(c\,x_i)\right) + a + e,
\end{equation}
где $a = 20$, $b = 0{,}2$, $c = 2\pi$. Минимум, равный нулю, достигается в начале
координат. Вдали от минимума функция почти плоская и плохо подсказывает
направление спуска, поэтому она проверяет, насколько точно метод выходит к
минимуму.

Функция Букина N.6 определена только для двух переменных:
\begin{equation}\label{eq:bukin6}
  f_{\text{buk}}(x_1, x_2) = 100\sqrt{\bigl|x_2 - 0{,}01\,x_1^{2}\bigr|}
  + 0{,}01\,\bigl|x_1 + 10\bigr|.
\end{equation}
Её минимум лежит вдоль узкого изогнутого оврага: малый шаг в сторону от него
резко увеличивает значение функции. Эта функция проверяет, насколько хорошо
алгоритм движется вдоль сложных по форме областей пространства решений.

Такой набор охватывает задачи разного характера --- от простой функции
с единственным минимумом до функции с множеством локальных минимумов и до
функции с узким оврагом, --- что позволяет всесторонне оценить поведение
метода.

\subsection{Алгоритм биогеографической оптимизации}
\label{subsec:bbo}

\subsection{Модификации модели миграции}
\label{subsec:migration}

\subsection{Модификации базового алгоритма}
\label{subsec:bbo-mods}

\subsection{Метод роя частиц и сравнение с биогеографической оптимизацией}
\label{subsec:pso}

\subsection{Обзор программ-аналогов}
\label{subsec:analogs}

% =====================================================================
%  2. КОНСТРУКТОРСКАЯ ЧАСТЬ
% =====================================================================
\section{Конструкторская часть}\label{sec:design}

\subsection{Функциональные требования к программному комплексу}
\label{subsec:requirements}

\subsection{Архитектура программного комплекса}
\label{subsec:architecture}

\subsection{Проектирование протокола обмена данными}
\label{subsec:protocol-design}

\subsection{Модель хранения данных}
\label{subsec:data-model}

\subsection{Проектирование пользовательского интерфейса}
\label{subsec:ui-design}

% =====================================================================
%  3. ТЕХНОЛОГИЧЕСКАЯ ЧАСТЬ
% =====================================================================
\section{Технологическая часть}\label{sec:technology}

\subsection{Обоснование выбора технологий}
\label{subsec:stack}

\subsection{Реализация вычислительного ядра сервера}
\label{subsec:core}

\subsection{Реализация очереди задач и обработчика}
\label{subsec:queue}

\subsection{Реализация алгоритмов оптимизации}
\label{subsec:algo-impl}

\subsection{Реализация мобильного клиента}
\label{subsec:client-impl}

\subsection{Реализация визуализации результатов}
\label{subsec:visualization}

% =====================================================================
%  4. АНАЛИТИЧЕСКАЯ ЧАСТЬ
% =====================================================================
\section{Аналитическая часть}\label{sec:analysis}

\subsection{Тестирование программного обеспечения}
\label{subsec:testing}

\subsection{Вычислительные эксперименты и анализ результатов}
\label{subsec:experiments}

% =====================================================================
%  ЗАКЛЮЧЕНИЕ — ~1 страница, не нумеруется
% =====================================================================
\structelem{ЗАКЛЮЧЕНИЕ}
% TODO: краткие выводы по каждой задаче с упором на результат,
%       оценка полноты выполнения задания, перспективы развития.

% =====================================================================
%  СПИСОК ИСПОЛЬЗОВАННЫХ ИСТОЧНИКОВ — 15-20 наименований, ГОСТ 7.32
% =====================================================================
\clearpage
\phantomsection
\addcontentsline{toc}{section}{СПИСОК ИСПОЛЬЗОВАННЫХ ИСТОЧНИКОВ}
\begin{center}\large\bfseries СПИСОК ИСПОЛЬЗОВАННЫХ ИСТОЧНИКОВ\end{center}
\vspace{8pt}
\printbibliography[heading=none]

% =====================================================================
%  ПРИЛОЖЕНИЯ
% =====================================================================
\structelem{ПРИЛОЖЕНИЯ}
% TODO: листинги ключевых модулей, руководство пользователя (при
%       необходимости), графический материал.

\end{document}